% Pursuing Stacks (cotes 134-2 to 134-8): archive cote 134-2, site page 36 (PDF page 37).
% letter to Quillen. Draft compared: ps1.tex:1096-1169.
% AI-transcribed by Claude Opus 5.5 (claude-opus-5-5, adaptive thinking, effort high), Message Batches API, batch ps1 (msgbatch_01Rv5hCaNTtUt4JHx1yqoHPL), 2026-11-12; not yet checked by eye.
\scanpage{134-2/36}
\margin{12}
\mspage{12}

\inserted{implies} inductively that an automorphism of the system of $\underline{F}_i$'s  related by the
source and target maps $\underline{s}_1^i$ , $\underline{t}_1^i$  is the identity.  Thus it is correct afte\unsure{r}
% NOTE: the right edge of the scan cuts off the end of the line above
all, it seems,  that the category structure of a coherator implies all
other structure relevant to us\hand{(21)}.
\margin{\unsure{Maybe f th} it is safer to give moreover the $\underline{C}_o$ subcategory}

I do believe that the description given so far of what I mean by a
coherator, \xout{\illegible{1}} namely something acting like a kind of pattern in
order to define a corresponding notion of ``stacks'' (which in turn should
be the basic coefficient objects in non commutative homological algebra,
as well as a convenient description of homotopy types) embodies some of
the essential features of the theory \inserted[typed]{still in embryo} that wants to be developped. It is
quite possible of course that some features are lacking still ,  for in-%
stance that some extra conditions have to be imposed upon $\underline{C}$ ,  possibly
of a very different nature from mere irredundancy conditions (which, I
feel, are kind of irrelevant in this set-up).  Only by pushing ahead
and working out at least in outline the main aspects of the formalism of
stacks,  will it become clear whether or not extra conditions on $\underline{C}$ are
needed.  I would like at least to make a commented list of these main
aspects,  and possibly do some heuristic pondering on some of these in
the stride,  or afterwards.  For today it seems a little late though  -
I have been pretty busy with non mathematical work most part of the day,
and the next two days I'll be busy at the university.  Thus I guess I'll
send off this letter tomorrow,  and send you \unsure{a\xout{n}|an} \inserted{later} elaboration  (presumably
much in the style of this unending letter)  if you are interested.  In
any case I would appreciate any comments you make  -  that's why I have
been writing you after all !  I will probably send copies to Ronnie
Brown,  Luc Illusie and Jean Giraud,  in case they should be interested
(I guess at least Ronnie Brown is).  Maybe the theory is going to take
off after all,  in the long last !

Very cordially yours \hand{Alexander}

PS  (\unsure{23.2|25.2})  I noticed a rather silly mistake in the notes of two days
ago,when stating that the categories $\underline{C}_n$ admit fiber\xout{ed} products \unsure{\over{-}{:}} what is
\margin{\unsure{(1')}}
true is that the category $\underline{C}_o$ has fiber\xout{ed} products (by construction \inserted[typed]{practically}) \hand{\unsure{(1')}}
and that these are fiber\xout{ed} products also in the categories $\underline{C}_n$ (by con-%
struction equally) i.e. that the inclusion functors $\underline{C}_o \to \underline{C}_n$ commute to
fiber\xout{ed} products.  Stacks in a category C correspond to functors
$\underline{C}_\infty \to C$  whose \uline{restriction to $\underline{C}_o$} commutes with fiber products. I carried
the mistake along in the \xout{\illegible{1}} yesterday notes  -  it does'nt really
change anything substantially.  I will have to come back anyhow upon the
basic notion of a coherator ...

\margin{36}
